\documentclass[11pt]{article}
\usepackage[margin=1in]{geometry}
\usepackage{textcomp}
\usepackage{amsmath,amssymb,amsthm,booktabs,hyperref}
\hypersetup{colorlinks=true,linkcolor=black,citecolor=black,urlcolor=blue}
\newtheorem{theorem}{Theorem}
\newtheorem{lemma}[theorem]{Lemma}
\newtheorem{proposition}[theorem]{Proposition}
\theoremstyle{remark}
\newtheorem{remark}[theorem]{Remark}
\newcommand{\R}{\mathbb{R}}
\newcommand{\M}{\mathfrak{M}}
\newcommand{\N}{\mathfrak{N}}
\newcommand{\s}{\mathfrak{s}}
\newcommand{\tf}{\mathfrak{t}}
\newcommand{\B}{\mathfrak{B}}
\newcommand{\D}{\mathfrak{D}}
\newcommand{\gr}{\nabla_{\!r}}
\DeclareMathOperator{\dv}{div}
\DeclareMathOperator{\curl}{curl}
\title{The P\'olya--Szeg\H{o} conjecture for convex pentagons\\[4pt]
\large Computer-assisted proof candidate, not peer reviewed\\
(Version 2.0, 2026-11-28)}
\author{Puyi Sun\thanks{ORCID 0009-0004-1644-5300. See Section~\ref{sec:ai} for the roles of AI systems in this work.}}
\date{}
\begin{document}
\maketitle

\begin{abstract}
For a planar domain $\Omega$ let $F(\Omega)=\lambda_1(\Omega)\,|\Omega|$, $\lambda_1$ the first Dirichlet eigenvalue of the
Laplacian. P\'olya and Szeg\H{o} conjectured that among polygons with $n$ sides the regular one minimises $F$; this is
classical for $n=3,4$; a recent preprint \cite{wang} claims a proof for all $n\ge5$ (for simple polygons with at most $n$
sides). Independently of it, we describe a computer-assisted proof candidate of the statement for \emph{convex}
pentagons: $F(P)\ge F(P_{\mathrm{reg}})$ for every convex pentagon $P$, with equality if and only if $P$ is similar to the regular
one. The proof has three parts: an explicit ball around the regular pentagon in a six-dimensional chart, where a
Lehmann--Goerisch bound with a trial family containing the true eigenfunction is exact at the centre, with vanishing
gradient; an annulus,
covered by boxes on each of which one trial family proves the bound for all pentagons of the box; and the rest, treated
in a chart in which a pentagon is a quadrilateral with a corner cut off. Every decision of the computations is an
inequality between exact rational numbers or a statement of ball arithmetic, and the certificates can be replayed box
by box. All certificates of the three parts were re-derived and replayed by independently written programs of a
second-party AI reviewer. Non-convex pentagons are not covered.
\end{abstract}

\section{Result and structure of the proof}
P\'olya and Szeg\H{o} proved that the equilateral triangle and the square minimise $F$ among triangles and
quadrilaterals and conjectured the same for $n$-gons, $n\ge5$ \cite{ps}. Bogosel and Bucur reduced the case of fixed
$n$ to a finite computation \cite{bb-jep} and proved, with interval arithmetic, that the regular pentagon and hexagon
are local minimisers \cite{bb-val} (we did not find an explicit size of the neighbourhood there); it is not known a priori that an optimal
polygon is convex \cite{bb-jep}. A public repository maintained under the GitHub account \texttt{xfliu} \cite{liu}
contains a computer-assisted proof of strict local minimality of the regular pentagon along every direction for
parameters $t\le1/6250$ (an explicit neighbourhood in its own coordinates), based on the one-vector Lehmann--Goerisch
bound. The conjecture is known for convex polygons with sufficiently many sides \cite{cghl}, and its local version for
the torsional rigidity for all $n\ge5$ \cite{bbf}; the same authors have since posted a preprint claiming an analytic proof of the
torsional (Saint-Venant) inequality for all $n\ge3$ \cite{bbf2}, in which they regard a transparent, human-verifiable
proof of the eigenvalue inequality as a goal that remains worthwhile. A preprint of H.~Wang \cite{wang} (17 November 2026) claims a proof of
the P\'olya--Szeg\H{o} inequalities for the first eigenvalue and the torsional rigidity for every $n\ge5$ and all simple
polygons with at most $n$ sides; it contains the case of convex pentagons treated here. It has
not been peer reviewed, and we have not checked it. The present work does not depend on it: it is an independent,
computer-assisted argument for the convex case by a different method, whose certificates can be replayed.
This work was carried out independently of \cite{wang}: an earlier stage of it on the same problem and in the same
chart (another certification method and certificates for a few boxes of pentagons) was released on 1 November 2026 as
version~1.0 of the accompanying repository \cite{v10}, and we learned of
\cite{wang} on 23 November 2026, after the certificates of Sections~\ref{sec:near}--\ref{sec:labels} had been
computed. No argument of \cite{wang} is used here, and we make no priority claim.

\begin{theorem}\label{thm:main}
For every convex pentagon $P$ (simple, of positive area): $F(P)\ge F(P_{\mathrm{reg}})$, with equality if and only if $P$
is similar to the regular pentagon.
Moreover $18.9191028942<F(P_{\mathrm{reg}})<18.9191047951$.
\end{theorem}

\paragraph{Chart.} $z\in\R^6\mapsto P(z)$ = the polygon with vertices $(0,0),(1,0),(z_1,z_2),(z_3,z_4),(z_5,z_6)$.
Let $c$ be the point of the regular pentagon and $C$ the point of binary64 numbers nearest to it, modified so that
$C_3=\frac12$, $C_5=1-C_1$, $C_6=C_2$ hold exactly; $|C-c|_\infty\le2^{-50}$. A strictly convex pentagon $P$ has ten
labellings (five cyclic ones, with and without reflection), each giving a chart point $z_h(P)$; put
$d(P)=\min_h|z_h(P)-C|_2$. Let $F^*$ be the exact terminating decimal $18.919134195733914$; it is larger than the upper
bound for $F(P(c))$ in Theorem~\ref{thm:L} (so $F(P_{\mathrm{reg}})\approx18.919104$ and the common comparison
threshold $F^*$ are different numbers). In the near producer, the binary64 literal nearest to this decimal is slightly
larger (by about $8.3\cdot10^{-17}$). Some far producer routines use an outward bound derived from the regular-pentagon
certificate and others use individual table targets. A certificate for any larger target implies the stated
inequality; stored table targets and binary64 geometric inputs are interpreted as their exact encoded rationals. The
independent re-certification uses the exact terminating decimal as the common target and verifies the stored table
bounds separately.

\begin{theorem}[local]\label{thm:L}
For every $z$ with $0<|z-c|_2<R=52429/2^{20}$ $(>\frac1{20})$, $P(z)$ is a simple pentagon and $F(P(z))>F(P(c))$. Moreover
$18.9191028942<F(P(c))<18.9191047951$.
\end{theorem}
\begin{theorem}[near]\label{thm:N}
For every $z$ with $\frac1{20}\le|z-C|_2<\frac35$, $P(z)$ is a simple pentagon, star-shaped with respect to an
interior point, and $F(P(z))>F^*$.
\end{theorem}
\begin{theorem}[far]\label{thm:F}
Every strictly convex pentagon $P$ satisfies $F(P)>F^*$ or $d(P)<\frac35$.
\end{theorem}
\begin{proof}[Proof of Theorem~\ref{thm:main} from Theorems~\ref{thm:L}--\ref{thm:F}]
A convex pentagon with a flat vertex is a quadrilateral or a triangle, and $F\ge2\pi^2>F^*>F(P(c))$ for those
\cite{ps}. Let $P$ be strictly convex. $F$ is invariant under similarities and reflections, so $F(P)=F(P(z))$ for each
of its chart points. If $d(P)\ge\frac35$, then $F(P)>F^*>F(P(c))$ by Theorem~\ref{thm:F}. Otherwise let
$z=z_h(P)$ with $|z-C|<\frac35$. If $|z-C|\ge\frac1{20}$, Theorem~\ref{thm:N} applies. If $|z-C|<\frac1{20}$, then
$|z-c|<\frac1{20}+\sqrt6\,2^{-50}<R$, and Theorem~\ref{thm:L} gives $F(P)>F(P(c))$ unless $z=c$.
\end{proof}

All three theorems rest on one inequality (Lemma~\ref{lem:LG}) applied to piecewise polynomial trial functions on a fan
of triangles; the certificates are positivity statements for polynomials with rational coefficients
(Section~\ref{sec:cert}). Theorem~\ref{thm:F} uses in addition classical bounds and a covering argument
(Sections~\ref{sec:cut}--\ref{sec:cover}).

\section{The lower bound}
\begin{lemma}[one-vector Lehmann--Goerisch bound]\label{lem:LG}
Let $\Omega\subset\R^2$ be a bounded domain with Dirichlet eigenvalues $0<\lambda_1\le\lambda_2\le\dots$, and
$\nu\le\lambda_2$. Let $u\in H^1_0(\Omega)$, $u\ne0$, and $w\in H(\dv;\Omega)$ with $-\dv w=u$, and
$A_0=\|\nabla u\|^2$, $A_1=\|u\|^2$, $A_2=\|w\|^2$. If $A_0<\nu A_1$, then $\nu A_2-A_1>0$ and
\[
\lambda_1\ \ge\ \frac{\nu A_1-A_0}{\nu A_2-A_1}.
\]
\end{lemma}
\begin{proof}
Let $\psi\in H^1_0$ solve $-\Delta\psi=u$. For $f\in H^1_0$: $\int\nabla f\cdot w=\int fu=\int\nabla f\cdot\nabla\psi$,
so $w-\nabla\psi\perp\nabla H^1_0$ and $A_2\ge\|\nabla\psi\|^2=:B_2$. With $u=\sum c_k\varphi_k$ in an orthonormal
eigenbasis and $p_k=c_k^2$: $A_1=\sum p_k$, $A_0=\sum\lambda_kp_k$, $B_2=\sum p_k/\lambda_k$. Put $D=A_0-\nu A_1=
\sum p_k(\lambda_k-\nu)<0$ and $G_0=A_0-2\nu A_1+\nu^2B_2=\sum p_k(\lambda_k-\nu)^2/\lambda_k$. Then $\lambda_1<\nu$, and
since $\lambda_k\ge\nu$ for $k\ge2$: $-D\le p_1(\nu-\lambda_1)$ and $\lambda_1G_0\ge p_1(\nu-\lambda_1)^2\ge
(-D)(\nu-\lambda_1)$, i.e.\ $\lambda_1(G_0-D)\ge\nu(-D)$ with $G_0-D\ge-D>0$. Finally $G:=A_0-2\nu A_1+\nu^2A_2\ge G_0$,
$t\mapsto\nu(-D)/(t-D)$ is decreasing, and $\nu(-D)/(G-D)=(\nu A_1-A_0)/(\nu A_2-A_1)$.
\end{proof}
This is the case of one trial vector of the method of Lehmann and Goerisch \cite{lehmann,gh,bg}. For $u=\varphi_1$,
$w=\nabla\varphi_1/\lambda_1$ the bound equals $\lambda_1$. We use $\nu=\kappa_0/|\Omega|$ with
$\kappa_0=36.33682907<2\pi j_{0,1}^2$, admissible by the Krahn--Szeg\H{o} inequality
$\lambda_2(\Omega)|\Omega|\ge2\pi j_{0,1}^2$ \cite{krahn,hong}.

\section{Reference coordinates and the fan flux}
Fix a rational point $O$ near the centre of $P(c)$. For $z$ in the cube $\{|z-c|_\infty\le W\}$, $W=13/256$, the fan
triangles $T_k(z)=(O,V_k(z),V_{k+1}(z))$, $k=0,\dots,4$, have $\det\nolimits_k(z)=e_1\times e_2>0$
($e_1=V_k-O$, $e_2=V_{k+1}-O$; proved on the whole cube, Section~\ref{sec:side}); since the cube is connected and
contains $c$, the triangles tile $P(z)$. On $T_k$ write $p=O+ae_1+be_2$ with $(a,b)$ in the reference triangle
$a,b\ge0$, $a+b\le1$; functions on $P(z)$ are functions of $(k,a,b)$, and $\gr=(\partial_a,\partial_b)$. With
$n_{11}=|e_1|^2$, $n_{12}=e_1\cdot e_2$, $n_{22}=|e_2|^2$, $\det=\det_k$ put
\[
\M_z=\frac1{\det}\begin{pmatrix}n_{22}&-n_{12}\\-n_{12}&n_{11}\end{pmatrix},\qquad
\N_z=\M_z^{-1}=\frac1{\det}\begin{pmatrix}n_{11}&n_{12}\\n_{12}&n_{22}\end{pmatrix},
\]
\[
b_z(v,v')=\sum_k\det\nolimits_k\!\int vv'\ \Big(=\int_{P(z)}vv'\Big),\qquad
a_z(v,v')=\sum_k\int\gr v\cdot\M_z\gr v'\ \Big(=\int_{P(z)}\nabla v\cdot\nabla v'\Big),
\]
all integrals over the reference triangle. A vector field $w$ on $P(z)$ is represented by $\s=\det\cdot$(its
components in the basis $e_1,e_2$); then $\int_{P(z)}w\cdot w'=c_z(\s,\s'):=\sum_k\int\s\cdot\N_z\s'$,
$\nabla v$ corresponds to $\s=\M_z\gr v$, and $\int_{P(z)}w\cdot\nabla f=\sum_k\int\s\cdot\gr f$.

\emph{Fan flux.} For $v$ continuous and piecewise smooth and $\chi$ continuous let $h=(2+E)^{-1}v$, $E=a\partial_a+
b\partial_b$, i.e.\ $h(s,\theta)=s^{-2}\int_0^st\,v(t,\theta)\,dt$ along the rays from $O$, and
\[
\s_z(v,\chi)=\big(-\det\nolimits_k\,a\,h+\partial_b\chi,\ -\det\nolimits_k\,b\,h-\partial_a\chi\big)
\quad\longleftrightarrow\quad w=-(p-O)\,h+\curl\chi .
\]
Then $-\dv w=(2+E)h=v$ in every triangle; across the spokes the normal component of $(p-O)h$ vanishes and that of
$\curl\chi$ is the tangential derivative of the continuous $\chi$. Hence $w\in H(\dv;P(z))$, $-\dv w=v$, and for
$f\in H^1_0$
\begin{equation}\label{eq:IP}
\sum_k\int\s_z(v,\chi)\cdot\gr f=b_z(v,f).
\end{equation}

\section{Certificates}\label{sec:cert}
The notation of the previous section ($\det_k$, $\M$, $\N$, $a$, $b$, $c$, the fan flux) is used below for any polygon
$P$ with vertices $V_k$ and any point $O$ whose fan determinants $\det_k=(V_k-O)\times(V_{k+1}-O)$ are positive.

\paragraph{One polygon.} Let $P$ be a simple polygon and $O$ a point with all fan determinants $\det_k>0$ and winding
number one about $O$ (for instance a strictly convex polygon and an interior point); then the fan triangles tile $P$. Take $u$ continuous,
a polynomial of degree $N$ on every fan triangle (Bernstein form, coefficients $2^{-s}\times$ integers, zero on the
boundary), and $\chi$ of the same kind of degree $N+1$ without boundary condition. The pair $(u,\s(u,\chi))$ is
admissible in Lemma~\ref{lem:LG}, and $A_0,A_1,A_2$ are sums over the triangles of fixed rational quadratic forms of the
coefficient lists (they depend on $N$ only), with the four numbers $n_{11},n_{12},n_{22},\det$ of the triangle as
weights. With $\nu=\kappa_0/|P|$ the lemma gives a rational lower bound of $F(P)$. The coefficients come from a
floating-point Rayleigh--Ritz computation and a floating-point least-squares problem for $\chi$; the proof does not
depend on their quality.

\paragraph{A box of polygons.} Let the vertices be polynomial functions of parameters $\zeta$ in a box
$X=\{|\zeta_i-m_i|\le\rho_i\}$, $O$ fixed, and let the coefficient lists be affine in $\zeta$:
$u(\zeta)=u^0+\sum_i(\zeta_i-m_i)u^i$, likewise $\chi$ (the lists $u^i,\chi^i$ are numerical derivatives of the data at
the centre). If $\det_k(\zeta)>0$ on $X$ for all $k$, the winding number of the fan about $O$ is constant on $X$; it is
$1$ as soon as one $P(\zeta)$ is a simple polygon (a pentagram, with winding number $2$, also has positive fan
determinants). Then every $P(\zeta)$, $\zeta\in X$, is a simple polygon, star-shaped about $O$, the triangles tile it, and
the transplanted pair is admissible on it. With $\D=\prod_k\det_k$, $|P|=\frac12\sum_k\det_k$,
\[
P_D:=\D\,(\kappa_0A_1-|P|A_0),\qquad P_C(L):=\D\,\big(|P|(\kappa_0A_1-|P|A_0)-L(\kappa_0A_2-|P|A_1)\big)
\]
are polynomials in $\zeta-m$ with rational coefficients, and by Lemma~\ref{lem:LG}:
\begin{equation}\label{eq:box}
\det\nolimits_k>0\ (\forall k),\quad P_D>0,\quad P_C(L)>0\quad\text{on }X\qquad\Longrightarrow\qquad F(P(\zeta))>L\ \text{ for all }\zeta\in X ,
\end{equation}
provided that some $P(\zeta_0)$, $\zeta_0\in X$, is a simple polygon (so that the winding number is $1$); in
particular $F(P(\zeta))>L$ for every member of $X$ that is itself a simple polygon whose fan about $O$ has winding
number $1$, for instance a strictly convex polygon with $O$ in its interior.
\begin{lemma}[positivity on a box]\label{lem:pos}
Let $p(x)=p_0+g\cdot x+\frac12x^THx+\sum_{|\alpha|\ge3}p_\alpha x^\alpha$ and $\mu\in\R^n_{\ge0}$ with
$A=H+2\,\mathrm{diag}(\mu)$ positive definite. Then for $|x_i|\le\rho_i$:
$p(x)\ge p_0-\frac12g^TA^{-1}g-\sum_i\mu_i\rho_i^2-\sum_{|\alpha|\ge3}|p_\alpha|\rho^\alpha$.
\end{lemma}
\begin{proof}
$g\cdot x+\frac12x^THx=g\cdot x+\frac12x^TAx-\sum\mu_ix_i^2\ge-\frac12g^TA^{-1}g-\sum\mu_i\rho_i^2$.
\end{proof}
The multipliers $\mu$ come from floating point; positive definiteness is verified by an exact $LDL^T$ factorisation.
If the test fails, the box is bisected and the polynomials are re-expanded at the new centres (exact Taylor shift).

\section*{Part I: the neighbourhood of the regular pentagon (Theorem~\ref{thm:L})}
In this part $P(z)$ is the pentagon chart, $d=z-c$, and the fan has the five vertices of $P(z)$.

\section{The family: exact at the regular pentagon}\label{sec:family}
Let $u_0$ be a first eigenfunction of $P(c)$, $\lambda_0=\lambda_1(P(c))$, $h_0=(2+E)^{-1}u_0$, and $\chi_0\in H^1$
with $\curl\chi_0=\nabla u_0/\lambda_0+(p-O)h_0$ on $P(c)$ (the right side is a divergence-free $L^2$ field on a
simply connected Lipschitz domain); equivalently $\s_c(u_0,\chi_0)=\M_c\gr u_0/\lambda_0$. Let $U_1,\dots,U_6$ be
continuous piecewise polynomials on the fan vanishing on the boundary, and $X_1,\dots,X_6$ continuous piecewise
polynomials. With $d=z-c$,
\[
v_z=u_0+\sum_id_iU_i,\qquad\chi_z=\chi_0+\sum_id_iX_i,\qquad\s_z=\s_z(v_z,\chi_z).
\]
Transplanted to $P(z)$, $(v_z,\s_z)$ is admissible in Lemma~\ref{lem:LG}. Let $|P|=|P(z)|$, $F_0=\lambda_0|P(c)|$,
\begin{align*}
\B_z\big((v,\s),(v',\s')\big)&=|P|(\kappa_0+F_0)\,b_z(v,v')-|P|^2a_z(v,v')-F_0\kappa_0\,c_z(\s,\s'),\\
\Psi(z)&=\B_z\big((v_z,\s_z),(v_z,\s_z)\big).
\end{align*}
If $\kappa_0b_z(v_z,v_z)>|P|a_z(v_z,v_z)$, Lemma~\ref{lem:LG} with $\nu=\kappa_0/|P|$ gives
\begin{equation}\label{eq:main}
F(P(z))-F_0\ \ge\ \frac{\Psi(z)}{\kappa_0c_z(\s_z,\s_z)-|P|\,b_z(v_z,v_z)},\qquad\text{denominator}>0 .
\end{equation}
At $z=c$: $a_c=\lambda_0b_c$ and $c_c(\s_c,\s_c)=b_c/\lambda_0$, so $\Psi(c)=b_c\,[\,|P|(\kappa_0+F_0)-|P|^2\lambda_0-
F_0\kappa_0/\lambda_0]=0$.

\begin{lemma}\label{lem:S}
$\nabla\Psi(c)=0$.
\end{lemma}
\begin{proof}
Let $L(z)=F_0+\Psi(z)/(\kappa_0c_z(\s_z,\s_z)-|P|b_z(v_z,v_z))$ and $U(z)=|P(z)|\,a_z(u_0,u_0)/b_z(u_0,u_0)$, the Rayleigh
quotient of the transplanted $u_0$. Both are rational functions of $z$ with fixed integrals as coefficients, smooth
near $c$, and $L\le F\le U$ near $c$ by \eqref{eq:main} and Rayleigh's principle, with $L(c)=U(c)=F_0$ (the hypothesis
of \eqref{eq:main} holds at $c$ because $\kappa_0>F_0$, hence in a neighbourhood). Let
$\varphi_g$, $g\in D_5$, be the ten relabelling maps of the chart: they are smooth near $c$, $\varphi_g(c)=c$ and
$F\circ\varphi_g=F$. Hence $L\le F\circ\varphi_g\le U\circ\varphi_g$ for every $g$, and averaging,
$L(z)-F_0\le\nabla U(c)\,M\,(z-c)+O(|z-c|^2)$ with $M=\frac1{10}\sum_gD\varphi_g(c)$. The matrices $D\varphi_g(c)$ form a
representation of $D_5$ and $M$ is the projection onto its fixed vectors; a computation in ball arithmetic gives
$\|M\|_\infty<10^{-100}<1$, hence $M=0$. (Alternatively: the tangent representation is $\rho_1\oplus2\rho_2$, without
trivial component.) So $\nabla L(c)\cdot(z-c)\le O(|z-c|^2)$ for all $z$, i.e.\ $\nabla L(c)=0$, and
$\nabla\Psi(c)=(\text{denominator at }c)\,\nabla L(c)=0$.
\end{proof}
No differentiability of $F$ is used. Consequently
\[
Q(d):=\Big(\prod_k\det\nolimits_k(c+d)\Big)\,\Psi(c+d)
\]
is a polynomial in $d$ \emph{without constant and linear terms}, and by \eqref{eq:main} Theorem~\ref{thm:L} follows
from: (i) $\det_k>0$; (ii) $\kappa_0b_z(v_z,v_z)>|P|a_z(v_z,v_z)$; (iii) $Q(d)>0$; for $0<|d|_2<R$.

\section{The residual form}
Let $\tilde u_0,\tilde\chi_0$ be computed piecewise polynomial approximations of $u_0,\chi_0$,
$\tilde v_z=\tilde u_0+\sum d_iU_i$, $\tilde\chi_z=\tilde\chi_0+\sum d_iX_i$, $\tilde\s_z=\s_z(\tilde v_z,\tilde\chi_z)$,
$\tilde\Psi(z)=\B_z$ of the computed pair with itself. Let $e=\tilde u_0-u_0\in H^1_0$, $e_h=(2+E)^{-1}e$,
$\hat\lambda(z)=F_0/|P(z)|$ and
\[
\rho_z=\N_z\tilde\s_z-\gr\tilde v_z/\hat\lambda(z)
\]
(the pull-back of $w-\nabla u/\hat\lambda$ for the computed pair; $\rho_c$ is the residual at $c$).
\begin{proposition}\label{prop:R}
$\displaystyle\Psi(z)=K(z)+\sum_k\int\gr e\cdot\Lambda_z+\sum_k\int e_h\,(a,b)\cdot\Theta_z+q_z(e,e)$, where
\begin{align*}
K&=\tilde\Psi+2F_0\kappa_0\sum_k\!\int\rho_z\cdot\M_c\rho_c-F_0\kappa_0\sum_k\!\int(\M_c\rho_c)\cdot\N_z(\M_c\rho_c),\\
\Lambda_z&=2F_0\big[(\kappa_0/\lambda_0)\,\M_c\rho_z-|P(z)|\,\M_z\rho_z-(\kappa_0/\lambda_0)\,\M_c\N_z\M_c\rho_c\big],\\
\Theta_z&=-2F_0\kappa_0\,(\det\nolimits_k(z)-\det\nolimits_k(c))\,\big[\rho_z-\N_z\M_c\rho_c\big],\\
q_z(e,e)&=|P|(\kappa_0+F_0)\,b_z(e,e)-|P|^2a_z(e,e)-F_0\kappa_0\,c_z(\tf',\tf'),\\
\tf'&=\M_c\gr e/\lambda_0-(\det\nolimits_k(z)-\det\nolimits_k(c))\,(a,b)\,e_h .
\end{align*}
\end{proposition}
\begin{proof}
The true pair is the computed pair minus $(e,\tf_z)$, $\tf_z=\tilde\s_z-\s_z=\s_z(e,e_\chi)$ with
$e_\chi=\tilde\chi_0-\chi_0$ (the terms with $U_i,X_i$ cancel). At $c$: $\tf_c=\tilde\s_c-\M_c\gr u_0/\lambda_0=
\M_c(\rho_c+\gr e/\lambda_0)$, and $\tf_z=\tf_c-(\det_k(z)-\det_k(c))(a,b)e_h$. Thus
$\Psi=\tilde\Psi-2\B_z((\tilde v,\tilde\s),(e,\tf_z))+\B_z((e,\tf_z),(e,\tf_z))$. By the definition of $\rho_z$,
$c_z(\tilde\s,\tf_z)=\sum\!\int\rho_z\cdot\tf_z+\hat\lambda^{-1}\sum\!\int\gr\tilde v\cdot\tf_z$, and
$\sum\!\int\gr\tilde v\cdot\tf_z=b_z(\tilde v,e)$ by \eqref{eq:IP} for the pair $(e,e_\chi)$. Also
$\gr\tilde v=\hat\lambda(\N_z\tilde\s-\rho_z)$, so $a_z(\tilde v,e)=\hat\lambda\sum\!\int\tilde\s\cdot\gr e-
\hat\lambda\sum\!\int\rho_z\cdot\M_z\gr e=\hat\lambda b_z(\tilde v,e)-\hat\lambda\sum\!\int\rho_z\cdot\M_z\gr e$ by
\eqref{eq:IP} for the computed pair. With $|P|^2\hat\lambda=|P|F_0$ and $F_0\kappa_0/\hat\lambda=\kappa_0|P|$ the four
terms with $b_z(\tilde v,e)$ cancel:
\[
\B_z\big((\tilde v,\tilde\s),(e,\tf_z)\big)=F_0|P|\sum_k\!\int\rho_z\cdot\M_z\gr e-F_0\kappa_0\sum_k\!\int\rho_z\cdot\tf_z .
\]
Insert $\tf_z=\M_c\rho_c+\tf'$, expand $c_z(\tf_z,\tf_z)$, and collect the terms without $e$, linear and quadratic in $e$.
\end{proof}
The point of the proposition: a term-by-term estimate of the error of $\Psi$ is useless (the individual terms are of
size $(\kappa_0+F_0)|P|A_1$ and the smallest eigenvalue of the quadratic form to be protected is $0.52$ in units of
$F$), whereas here $e$ is paired only with fields built from $\rho$, which are small.

\section{Bounds for the unknown function}
Normalise $u_0$ as the $L^2(P(c))$ projection of $\tilde u_0$ on the first eigenspace. Then $e\perp u_0$,
$\|\nabla e\|^2-\lambda_0\|e\|^2=A_0-\lambda_0A_1$ ($A_i$: the forms of the computed pair at $c$) and
$\|\nabla e\|^2\ge\lambda_2\|e\|^2\ge\nu\|e\|^2$, so with $\lambda_{lo}\le\lambda_0\le\lambda_{up}=A_0/A_1$
($\lambda_{lo}$ from Lemma~\ref{lem:LG} at $c$)
\[
\varepsilon_1^2:=\|\nabla e\|^2_{L^2(P(c))}\le\frac{A_0-\lambda_{lo}A_1}{1-\lambda_{up}/\nu},\qquad
\varepsilon_0^2:=\|e\|^2\le\varepsilon_1^2/\nu .
\]
On every reference triangle $\int(a+b)^2e_h^2\le\frac14\int e^2$ (in the coordinates $s=a+b$, $\theta=a/s$:
$da\,db=s\,ds\,d\theta$, $e_h=s^{-2}\int_0^st\,e\,dt$, Cauchy--Schwarz). With degree $16$:
$\varepsilon_1\le1.02\cdot10^{-3}$ ($4.6\cdot10^{-4}$ of $\|\nabla\tilde u_0\|$), $\varepsilon_0\le2.3\cdot10^{-4}$.

\section{Enclosure of the coefficients}
Everything is multiplied by $\D(z)=\prod_k\det_k(z)$ and expanded in $d'=z-C_q$, $C_q$ a rational point with
$|c-C_q|_\infty\le2^{-50}$. Polynomial arithmetic is exact (FLINT \cite{flint}); quantities that are irrational or not
known exactly are carried as polynomials with a coefficientwise radius: $F_0$ and $\lambda_0$ (with rigorous interval enclosures recorded in \texttt{constants.json}; the enclosure of $F_0$
has width about $1.9\cdot10^{-6}$), the
geometry of $P(c)$, the weights $d'_i-(c-C_q)_i$ of the lists, and $\rho_c$. The coefficient fields of $\D\Lambda_z$
and $\D\Theta_z$ lie in the span of the $6\cdot7$ functions $a\,h_b,\ b\,h_b,\ \partial_a\chi_b,\ \partial_b\chi_b,\
\partial_aU_b,\ \partial_bU_b$ ($b=0,\dots,6$) per triangle, whose Gram matrices are exact rationals. For the
coefficient of $d'^\alpha$,
\[
\Big|\sum_k\!\int\gr e\cdot\Lambda_\alpha+\sum_k\!\int e_h(a,b)\cdot\Theta_\alpha\Big|\le
\varepsilon_1\Big(\sum_k\!\int\Lambda_\alpha\cdot\N_c\Lambda_\alpha\Big)^{1/2}+
\frac{\varepsilon_0}2\Big(\sum_k\!\int\frac{|\Theta_\alpha|^2}{\det_k(c)}\Big)^{1/2},
\]
computed exactly for $|\alpha|\le3$ and replaced by a cruder majorant for higher degrees; $q_z(e,e)$ is bounded by a
majorant with $\varepsilon_0^2,\varepsilon_1^2$; the Taylor shift from $C_q$ to $c$ changes a coefficient by at most
$(2^{-50}/2^{-10})\,[B(d+2^{-10}\mathbf 1)-B(d)]_\alpha$ for any coefficientwise majorant $B$. The result is
$Q(d)=\sum_{|\alpha|\ge2}s_\alpha d^\alpha$ with $s_\alpha\in[T_\alpha-E_\alpha,T_\alpha+E_\alpha]$, $T_\alpha,E_\alpha$
rational ($12{,}968$ and $13{,}024$ numbers). The computed coefficients of degree $0$ and $1$ are discarded: the true ones
vanish by Section~\ref{sec:family}. (They lie within their error bounds, as they must.)

\section{Positivity}
\emph{Inner cube $|d_i|\le\rho=1/128$.} Write $Q(d)=d^T[\frac12H+\sum_id_iT_i]d+\text{rest}$ with $H$, $T$ from
$T_\alpha$, $|\alpha|=2,3$, and $|\text{rest}|\le\tau|d|^2$, $\tau=\tau_2+\tau_3+\tau_4$:
$\tau_2=\max_i[E_{2e_i}+\frac12\sum_{j\ne i}E_{e_i+e_j}]$, $\tau_3=\rho\sum_{|\alpha|=3}E_\alpha w_\alpha$,
$\tau_4=\sum_{|\alpha|\ge4}\rho^{|\alpha|-2}(|T_\alpha|+E_\alpha)w_\alpha$ ($w_\alpha=1$ for pure powers, $\frac12$
otherwise). The matrix $\frac12H+\sum d_iT_i-\tau I$ is affine in $d$, so it is positive definite on the cube if it is
at the $64$ vertices (exact $LDL^T$). In units of $F$ (that is, $Q$ divided by its normalising factor
$\D(c)\,(\kappa_0c_c(\s_c,\s_c)-|P(c)|\,b_c(v_c,v_c))\approx0.109$; the certificate stores $Q$ itself): the eigenvalues
of $H$ are $0.525,\,0.557,\,5.03,\,5.70,\,12.76,\,15.31$; $\tau=0.165$ ($0.057+0.013+0.095$); the smallest pivot is
$0.32$ (in the units of the certificate: $\tau=0.01806$, smallest pivot $0.03492$).

\emph{Shell.} The half cube $\{|d_i|\le W,\ d_3\ge0\}$ is bisected into $7{,}939$ boxes $m\pm r$: $32$ inside the
inner cube, $965$ with $\sum_i\max(|m_i|-r_i,0)^2\ge R^2$, and $6{,}942$ on which
$p(x)=\sum_{2\le|\alpha|\le6}T_\alpha(m+x)^\alpha$ satisfies $\min_{|x_i|\le r_i}p>\sum_\alpha E'_\alpha\prod_i
(|m_i|+r_i)^{\alpha_i}$ ($E'=E$ plus $|T_\alpha|$ for $|\alpha|>6$), proved by: constant term, exact Lagrangian bound of
the linear and quadratic terms ($-\frac12g^TA^{-1}g-\sum\mu_ir_i^2$ with multipliers $\mu_i\ge0$ and $A=H_p+2\,\mathrm{diag}(\mu)$ positive
definite; both are checked),
and the sum of the absolute values of the higher terms. All of this is exact rational arithmetic. The other half
follows from the reflection $(z_1,\dots,z_6)\mapsto(1-z_5,z_6,1-z_3,z_4,1-z_1,z_2)$, an isometry of the chart that fixes
$c$ and preserves $F$.

\section{Side conditions}\label{sec:side}
On the cube of half-width $W$: $\det_k>0$ by an exact polynomial test; and
$\kappa_0b_z(v)-|P|a_z(v)\ge(\kappa_0b_z(\tilde v)-|P|a_z(\tilde v))-2\kappa_0(b_z(\tilde v)b_z(e))^{1/2}-
|P|(2(a_z(\tilde v)a_z(e))^{1/2}+a_z(e))$ with $b_z(e)\le2\varepsilon_0^2$, $a_z(e)\le2\varepsilon_1^2$
($\det_k(z)\le2\det_k(c)$, $\M_z\preceq2\M_c$ on the cube), which is proved positive on the whole cube by an exact polynomial test (the allowance is less than $0.032$; at the
centre the main term is about $7.776$).

\section*{Part II: the annulus (Theorem~\ref{thm:N})}
\section{The near cover}\label{sec:near}
The map $m(z)=(1-z_5,z_6,1-z_3,z_4,1-z_1,z_2)$ is an isometry of $\R^6$ with $m(C)=C$ exactly, and $P(m(z))$ is the
mirror image of $P(z)$ with the vertices relabelled; so it suffices to prove the theorem for $z_3\ge\frac12$. Let $K_0$
be the cube with centre $C$ rounded to multiples of $2^{-10}$ and half-width $\frac58$; it contains the ball
$|z-C|<\frac35$, and the plane $z_3=\frac12$ is a plane of its dyadic grid. $K_0$ is bisected into boxes with dyadic
ends; a box $[\mathrm{lo},\mathrm{hi}]$ becomes a leaf with the first applicable label:
\emph{mirror}, if $\mathrm{hi}_3\le\frac12$; \emph{far}, if $\inf_{\rm box}|z-C|^2\ge\frac9{25}$; \emph{inner}, if
$\sup_{\rm box}|z-C|^2<\frac1{400}$; \emph{LG}, if \eqref{eq:box} holds on the box with $L=F^*$ for the vertex map
$\zeta=z$ (the vertices are the parameters; $P_D$ and $P_C$ have degree $13$ and about $13{,}000$ terms). The infimum
and supremum of $|z-C|^2$ over a box are exact rational numbers.

\begin{proof}[Proof of Theorem~\ref{thm:N} from a complete cover]
Let $\frac1{20}\le|z-C|<\frac35$ and, after applying $m$ if necessary, $z_3\ge\frac12$. Then $z$ lies in a closed leaf
with $\mathrm{lo}_3\ge\frac12$; such a leaf is not \emph{mirror}, and it is neither \emph{far} nor \emph{inner} because
it contains $z$. So it is an LG leaf. $P(z)$ is simple: in the regular pentagon $P(c)$ (unit edges) two distinct
vertices, two non-adjacent edges, and a vertex and an edge not containing it are at distance at least $1$ (the vertex
angles are obtuse, and the distance to the opposite edge exceeds $1$). Moving the endpoints of a segment by at most $a$
moves the segment by at most $a$ in Hausdorff distance. In each of the three cases the two objects use disjoint sets of
the moving vertices $V_2,V_3,V_4$, whose displacements $a,b$ satisfy $a+b\le\sqrt2\,|z-c|<0.85<1$, so the three
distances stay positive. Two adjacent edges meeting in a further point would overlap on their common line, which puts
an endpoint of one of them on the other, at distance $0$ from an edge not containing it; this is excluded. The fan
determinants are positive on the whole leaf, so the winding number about the fan centre $O$ is constant on the leaf; at
$z$ it is $1$ because $P(z)$ is simple. So $P(z)$ is star-shaped about the interior point $O$, the fan triangles tile it,
and \eqref{eq:box} gives $F(P(z))>F^*$. (Nothing is claimed for members of the leaf outside the annulus.)
\end{proof}
The cover has $51{,}434$ leaves ($38{,}280$ LG, $13{,}118$ far, $36$ inner) in $1{,}344$ stored blocks of side
$\frac5{16}$; the half-widths of the LG leaves are $20$--$40\,\%$ of the distance to $C$.

\section*{Part III: far from the regular pentagon (Theorem~\ref{thm:F})}
\section{The cut chart}\label{sec:cut}
Let $P$ be a strictly convex pentagon, counter-clockwise, $e_i=V_{i+1}-V_i$, and for an edge $k$:
$A_k=e_{k-1}\times e_k>0$, $B_k=e_k\times e_{k+1}>0$, $C_k=e_{k-1}\times e_{k+1}$. If $C_k>0$, the lines of the edges $k\mp1$
meet beyond the edge $k$ in a point $W$; the triangle $(V_k,W,V_{k+1})$ (the \emph{ear} of $k$) has area
$\mathrm{ear}_k=A_kB_k/(2C_k)$, and $Q_k=P\cup\mathrm{ear}$ is a strictly convex quadrilateral. Some edge has $C_k>0$: the
exterior angles sum to $2\pi$, so two consecutive ones sum to less than $\pi$. Since $P\subset Q_k$,
\begin{equation}\label{eq:F1}
F(P)\ \ge\ F(Q_k)\,|P|/|Q_k| .
\end{equation}
Call $k$ \emph{kept} if $C_k>0$ and $\mathrm{ear}_k\le\mathrm{ear}_{k'}$ for all $k'$ with $C_{k'}>0$; every $P$ has a kept
edge.

\paragraph{Coordinates.} Normalise $Q=Q_k$ by a similarity, and a reflection if necessary, to
$Q(q)=(q_0,q_1,q_2,q_3)$, $q_0=0$, $q_1=(1,0)$, $q=(q_2,q_3)\in\R^4$, such that $q_0q_1$ is a longest edge and $W=q_j$ with
$j\in\{3,0\}$. With $t,u\in(0,1)$ the fractions of the two edges at $W$ that are cut away,
\[
j=3:\ P=[q_0,\,q_1,\,q_2,\,q_3+t(q_2-q_3),\,(1-u)q_3],\qquad j=0:\ P=[(t,0),\,q_1,\,q_2,\,q_3,\,u\,q_3],
\]
and $|P|=|Q|(1-tu\,\tau_j(q))$ with $\tau_j$ the area of the corner triangle at $q_j$ divided by $|Q|$. Call $q$
\emph{admissible} if $Q(q)$ is strictly convex and its edges other than the base have length $\le1$; then
$q\in\mathfrak B=[0,2]\times[0,1]\times[-1,1]\times[0,1]$. So every strictly convex pentagon is, up to similarity and
reflection, $P(q,t,u;j)$ with $q$ admissible, $t,u\in(0,1)$, $j\in\{3,0\}$ and the cut edge kept.

\paragraph{Invariants.} With the edge vectors $e_0=(1,0)$, $e_1=q_2-q_1$, $e_2=q_3-q_2$, $e_3=-q_3$ of $Q$:
$q_{3y}=e_3\times e_0$, $q_{2y}=e_0\times e_1$, $K=e_1\times e_2$, $D=e_2\times e_3$, $M=e_3\times e_1$, $L=e_0\times e_2$,
$2|Q|=q_{2y}+D$; $Q$ is strictly convex iff $q_{2y},K,D,q_{3y}>0$; $\tau_3=D/(q_{2y}+D)$, $\tau_0=q_{3y}/(q_{2y}+D)$.

\paragraph{Three charts and the corner principle.} std: $w=q$. E2: $w=(q_{3x},q_{3y},a,h)$, $q_2=q_1+a(q_3-q_1)+h(0,1)$,
so that $K=h(1-q_{3x})$. E3: $w=(q_{2x},q_{2y},b,h)$, $q_3=b\,q_2+h(0,1)$, so that $D=h\,q_{2x}$. Put $\zeta=(w,t,u)$. In
each chart every vertex of $P$ is multi-affine in $\zeta$, and each coordinate moves all vertices along one common
direction. Hence every cross product of two vertex differences is multi-affine in $\zeta$ (its second derivative in
$\zeta_m$ is $2\,\partial_mv\times\partial_mv'=0$), and so are the area and all linear functionals of the vertices.
A multi-affine function on a box attains its extreme values at the corners. All box criteria below evaluate such
functions at the $64$ (or $16$) corners, where the vertices are exact dyadic rationals.

\section{Leaves and their labels}\label{sec:labels}
\paragraph{Domains.} The std roots are the cells of side $\frac1{16}$ of $\mathfrak B$, except those where one of
$q_{2y},q_{3y},K,D$ is $\le0$ at all $16$ corners (no strictly convex $Q$) or where an edge other than the base is longer
than $1$ on the whole cell ($57{,}093$ roots). A $q$-box of the std chart is handed over to E2 if $q_{3y}\ge\frac14$ on
it and $K<\frac18(1-q_{3x})$ at its corners, to E3 if $q_{2x}\ge\frac14$ and $D<\frac18q_{2x}$. For an admissible strictly
convex $q$ of such a box, $h=K/(1-q_{3x})\in(0,\frac18)$ and $a(q_3-q_1)=e_1-h(0,1)$, $|q_3-q_1|\ge\frac14$, give
$|a|\le4.5$; resp.\ $h=D/q_{2x}\in(0,\frac18)$, $b=q_{3x}/q_{2x}$, $|b|\le4$. The E2 roots are the cells of
$[-1,1]\times[\frac14,1]\times[-4.5,4.5]\times[0,\frac18]$ that may contain a strictly convex $Q$ ($15{,}646$), the E3
roots those of $[\frac14,2]\times[0,1]\times[-4,4]\times[0,\frac18]$ ($25{,}660$). For every root $c$ and $j\in\{3,0\}$
the column $c\times[0,1]^2$ is bisected into leaves.

\paragraph{The quadrilateral bound.} For a $w$-box let $\mathrm{LB}$ be the largest of the following lower bounds of
$F(Q)$ over its strictly convex members: $2\pi^2$ \cite{ps}; for an edge $k$ of $Q$ with $C_k>0$ on the box,
$(4\pi^2/\sqrt3)/(1+\mathrm{ear}_k/|Q|)$ ($Q$ lies in the triangle of the other three edge lines, and
$F\ge4\pi^2/\sqrt3$ for triangles \cite{ps}); the strip bounds $\pi^2|Q|/w^2$ and $\pi^2/(2w)$, $w=\max(q_{2y},q_{3y})$;
Makai's bound $\pi^2\mathrm{Per}^2/(16|Q|)$ \cite{makai} (std chart); the fibre bound (Lemma~\ref{lem:OD}) of the convex
hull of the corner quadrilaterals times the smallest corner area; and, in the std chart, a table value of the root
cell: for the cell or each of its $16$ sub-cells with strictly convex members, a closed-form bound or a number $L$ for
which \eqref{eq:box} holds on the (sub-)cell for the quadrilateral chart.

\paragraph{Labels.} Each label is an assertion about the members of the box with $Q(w)$ strictly convex and
$t,u\in(0,1)$.
The criteria are proved in Appendix~\ref{app:far}.
\begin{itemize}\setlength\itemsep{1pt}
\item[V:] no member. An invariant $q_{2y},q_{3y},K$ or $D$ is $\le0$ at the $16$ corners.
\item[E2, E3:] (std) handed over, see above.
\item[T:] $F(P)>F^*$. The certified quadrilateral bound satisfies $\mathrm{LB}>F^*$, the upper bound of $\tau_j$ is positive, and
$\sup t\cdot\sup u<\theta_T:=(1-F^*/\mathrm{LB})/\sup\tau_j$; then
$F(P)\ge\mathrm{LB}\,(1-tu\tau_j)>F^*$ by \eqref{eq:F1}.
\item[X:] the cut edge is not kept: another edge $k'$ has $C_{k'}>0$ and $\mathrm{ear}_{k'}<\mathrm{ear}_{\rm cut}$ on the
box. For the two edges not touched by the cut this reads $tu>g(w)$, namely
($j=3$) $k'=0$: $M>0$, $g=q_{2y}q_{3y}/(DM)$; $k'=1$: $L>0$, $g=q_{2y}K/(LD)$;
($j=0$) $k'=1$: $L>0$, $g=q_{2y}K/(q_{3y}L)$; $k'=2$: $M<0$, $g=DK/(q_{3y}(-M))$.
For the two shortened edges it reads $C'>0$, $G<0$ with
($j=3$) $(C',G)=(u\,q_{3y}+t(q_{2y}-q_{3y}),\ q_{3y}(1-2u)-tu(q_{2y}-q_{3y}))$ and $(tK-uM,\ K(1-2t)+tuM)$;
($j=0$) $(uM+t\,q_{2y},\ q_{2y}(1-2t)-tuM)$ and $(uD+t(q_{2y}-q_{3y}),\ D(1-2u)-tu(q_{2y}-q_{3y}))$;
these are multi-affine in $\zeta$. The denominators of $g$ are certified positive on the box after the indicated sign
choices. In the chart E3 at the face $q_{2y}=0=h$ (where all five vertices can collapse to
one point) the first group is homogeneous in $(q_{2y},h)$ and is verified for all ratios $h/q_{2y}\in(0,\infty)$.
\item[TX:] (whole column) the strict separation $\theta_T>\sup g$ (with the positivity assumptions of T and of X): a member with
$tu<\theta_T$ has $F(P)>F^*$, a member with $tu\ge\theta_T>g$ is not kept.
\item[S:] $F(P)>F^*$ by a strip bound: (a) for a non-zero vector $n$ let $W_n=\max_{i,k}|\langle V_i-V_k,n\rangle|$, so
that $P$ lies in a strip of width $W_n/|n|$; if $|P|\ge\tau W_n^2/|n|^2$ on the box (certified for every vertex pair from
the corners, the function $|P|-\tau|n|^{-2}\langle V_i-V_k,n\rangle^2$ being concave in each coordinate) and
$\pi^2\tau>F^*$, then $F(P)\ge\pi^2\tau>F^*$ (Appendix~\ref{app:far}); (b) for an edge of length $l>0$ with $\max_kc_k>0$,
$F\ge\pi^2l^2/(2\max_kc_k)$, $c_k=(V_b-V_a)\times(V_k-V_a)$ ($P$ lies in a strip of width $H=\max_kc_k/l$ and contains a
triangle of area $lH/2$), and the label requires the verified lower bound to be strictly greater than $F^*$.
\item[R:] $F(P)>F^*$: for two non-adjacent edges $m,m+2$ with $C_m>0$ and $C_{m+2}>0$, $P$ lies in the triangle $T$ of the other three
edge lines, $|T|=|P|+\mathrm{ear}_m+\mathrm{ear}_{m+2}$, and $F(P)\ge(4\pi^2/\sqrt3)\,|P|/|T|$; the label requires this
verified lower bound to be strictly greater than $F^*$.
\item[O:] $F(P)>F^*$ by the fibre bound of the convex hull of the $320$ corner vertices (which contains $P(\zeta)$)
times the smallest corner area; the label requires this verified lower bound to be strictly greater than $F^*$.
\item[N:] $d(P)<\frac35$. For a labelling $h$ the chart coordinates are $u_p=N_p/D$ with $N_p=W_p-W_0$, $D=W_1-W_0$
complex multi-affine; with $r_p$ the value of $u_p$ at the centre, $|u_p-r_p|\le\rho_p:=\max_{\rm corners}|N_p-r_pD|/
\min_{\rm corners}\mathrm{Re}(D\bar e)$ for a unit vector $e$ (the denominator certified strictly positive), and the test
is $\sum_p(|r_p-C_p|+\rho_p)^2<\frac9{25}$ with outward rounding, $C_p$ the complex pair of $C$ that belongs to the
labelling.
\item[C:] $F(P)>F^*$: \eqref{eq:box} holds on the box with $L=F^*$ for the vertex map of the chart ($P_C$ has degree
up to $20$ and up to $41{,}000$ terms).
\end{itemize}

\begin{lemma}[fibre bound]\label{lem:OD}
Let $P$ be a bounded convex domain, $(a,b)$ its projection to the $x$-axis and $w(x)$ the length of its section at
$x$. Then $\lambda_1(P)\ge\mu:=\inf\{\int_a^b(f'^2+Vf^2):f\in C_c^\infty(a,b),\int f^2=1\}$ with $V=\pi^2/w^2$. If $V_t\le V$
is piecewise constant and there is $\varphi\in C^1[a,b]$, piecewise $C^2$, positive on $(a,b)$, with
$-\varphi''+V_t\varphi=m\varphi$ on every piece, then $\mu\ge m$.
\end{lemma}
\begin{proof}
For $u\in C_c^\infty(P)$, by Fubini and the one-dimensional Dirichlet inequality on every vertical section,
$\int|\nabla u|^2\ge\int\!\!\int(u_x^2+V(x)u^2)\,dx\,dy$; for fixed $y$ the horizontal slice of $u$, extended by zero,
lies in $C_c^\infty(a,b)$, so the inner integral is at least $\mu\int u^2\,dx$. For the second claim let $f\in
C_c^\infty(a,b)$ and $v=f/\varphi$, which has compact support in $(a,b)$ because $\varphi>0$ there; since $\varphi$ is
$C^1$ (values and derivatives match at the nodes), integration by parts on every piece leaves no boundary terms and
gives $\int(f'^2+V_tf^2)=\int\varphi^2v'^2+m\int f^2\ge m\int f^2$; and $V\ge V_t$.
\end{proof}
In the computation $w$ is bounded above on each of $128$ intervals by differences of edge lines, and $\varphi$ is
propagated by the explicit solution formulas in ball arithmetic; its positivity is certified on every open piece, not
only at the nodes (a hyperbolic or affine piece has at most one zero, so positive values at both ends suffice; a
trigonometric piece is in addition shorter than half a period). The independent re-certification uses its own routes
for these labels, listed with their conditions in the second party's \texttt{LABEL\_ROUTES.md}.

\section{The cover argument}\label{sec:cover}
\begin{proof}[Proof of Theorem~\ref{thm:F} from a complete cover]
Let $P$ be strictly convex with a kept edge and coordinates $(q,t,u;j)$, $q$ admissible. The cell of $\mathfrak B$
containing $q$ is a std root, so $(q,t,u)$ lies in a leaf of its column for $j$. If the leaf is E2 (E3), the edge-chart
coordinates of $q$ lie in the root range of that chart and in a root, and we pass to a leaf of that column. The leaf is
not V, because $Q(q)$ is strictly convex, and not X, because the cut edge is kept and X asserts a strictly smaller ear;
for the same reason the second alternative of TX is excluded. Hence the label is T, TX, S, R, O or C, and $F(P)>F^*$,
or it is N, and $d(P)<\frac35$.
\end{proof}

\section{Computations, verification and status}\label{sec:status}
\emph{(Status on 27 November 2026.)}
\begin{center}\small
\begin{tabular}{@{}lp{3.1cm}p{3.6cm}p{5.8cm}@{}}\toprule
part & computation & second computer & second-party review\\\midrule
Thm.~\ref{thm:L} & done ($6$ minutes) & yes, identical leaf list & passed in full: analysis; final positivity proof replayed independently; all $13{,}024$ coefficient enclosures recomputed independently and the positivity proof re-proved with them\\
Thm.~\ref{thm:N} & done & yes, self-check passed & certificate formulas, simplicity argument, mirror, skip tests and cover argument passed; structure of all blocks re-checked independently; an independently written exact certificate program re-proved all $38{,}280$ near targets and its stored witnesses were replayed; signed (PENT-26, 2026-11-28)\\
Thm.~\ref{thm:F} & done ($1{,}260$ blocks, $116$ CPU hours, no unresolved leaf) & $225$ blocks recomputed with matching results; root ranges, tiling and every closed-form label of the full run re-proved by a second implementation & Lemma~\ref{lem:LG}, \eqref{eq:box}, chart, kept representation, corner principle, ear comparisons and cover argument passed; the structure of the whole cover re-checked independently; an independently written exact program re-proved all $377{,}515$ certified boxes, $185{,}100$ table entries and $1{,}346{,}929$ labels, and its stored witnesses were replayed; signed (PENT-26, 2026-11-28)\\\bottomrule
\end{tabular}\end{center}
\noindent\textbf{The three computations are complete, have passed the producer-side checks, and were re-certified in full
by the second party's independently written programs; Theorem~\ref{thm:main} is signed by the second party (PENT-26, 2026-11-28).
It remains a computer-assisted proof candidate: it has not been peer reviewed, and the mathematics has not been
reviewed by human experts.}
The far cover consists of $1{,}724{,}444$ boxes: $377{,}515$ certified by \eqref{eq:box} and $1{,}346{,}929$ leaves with
closed-form labels. A second implementation written from the statement of the criteria re-proves every one of these
labels in exact rational arithmetic ($13$ of them only after one bisection of all six coordinates, because its
quadrilateral bounds are weaker than the producer's); the producer-side checks of the certified boxes and of the table
are described below. Producer-side replay of the certified boxes has the scope listed in the accompanying
manifest: the certified boxes of $6$ std blocks ($1{,}538$ boxes) were checked by sampling only. The independent full
re-certification is described below.

Producer-side checks (the same programs as the producer, run again; not a second implementation).
Theorem~\ref{thm:L}: a checker using the Python standard library only re-proves the final positivity statement from the
exported coefficients and boxes. Near: the stored blocks are re-enumerated, every block is re-tiled exactly, the skip
labels are re-tested, and all $38{,}280$ LG leaves of the $1{,}344$ blocks were re-proved from scratch with new trial
families. Far: root coverage and tiling were checked throughout (every other cell re-classified; exact rationals);
all closed-form labels were re-proved by a separate producer-side implementation; all table entries underwent the
stated producer checks. The producer-side re-proof of certified boxes was full on $422$ of the $428$ std blocks
($292{,}308$ boxes) and on $31$ E2 blocks ($1{,}871$ boxes); the other certified boxes were checked by sampling ($25$ per
block), as detailed in the run log. Independent full re-certification and exact replay are reported separately. Two floating-point tests that do
not use the label code: at $21{,}000$ random points of random leaves the assertion of the label holds; $20{,}000$ random
convex pentagons are reconstructed from their coordinates $(q,t,u;j)$, with $q$ in a root. End-to-end tests on the
complete covers (floating point): $110{,}000$ random strictly convex pentagons (among them $30{,}000$ adversarial shapes and
$30{,}000$ near the sphere $|z-C|=\frac35$) are followed through the kept representation into a leaf of the far cover
and, where the leaf hands over, of an edge chart; $100{,}000$ random points of the annulus of Theorem~\ref{thm:N} are
followed into a leaf of the near cover. At every point the assertion of the final leaf holds (a pointwise lower bound
of $F$ above $F^*$, or the chart distance below $\frac35$), and no point lies in a leaf of the wrong kind.

\paragraph{Independent re-certification of Parts~II and~III (22--27 November 2026).} The second party's programs,
written independently of the producer code, enumerate the $1{,}947{,}824$ targets of the near and far covers from the
frozen input archive (identified by its SHA-256 in the release), search a witness of their own for each target, and a
separate replay program re-checks every stored witness in exact arithmetic. They were run on the author's computer from
22 November 2026. In the main run every target was certified except $13$ near LG boxes ($38{,}267$ of the $38{,}280$
near targets). One of these $13$ was certified in a first supplementary round with the trial degree raised by two, and
the other twelve in a second supplementary round restricted to them, which uses the following form of witness. Each
original LG box is either certified in full or partitioned exactly into closed dyadic subboxes. A near-only
supplementary subbox may be labelled \emph{outside} when its exact squared distance from $C$ is at least
$(3/5+2^{-40})^2$. Such a subbox cannot contain the point of the annulus under consideration. Every remaining point of
the annulus is contained in an LG subbox. At that point the polygon is simple by Lemma~\ref{lem:S}; positive fan
determinants throughout that subbox make the winding number about the fan centre constant and equal to one. The LG
inequality therefore applies at the point. No bound is asserted on an outside subbox. (The producer's certificates of
Section~\ref{sec:near} hold on the whole of each box; outside subboxes occur only in this re-certification, and the
proof of Theorem~\ref{thm:N} does not use the far cover.) The tree of the first supplementary round has $11$ LG leaves; the twelve trees of the second
round consist of $86$ LG subboxes and $109$ outside subboxes. The replay is complete in all nine stages: each of the
$1{,}947{,}824$ targets replayed as certified ($299{,}994$ labels of the standard chart through the quadrilateral table,
which is itself certified in full), with no rejected witness and no unscanned shard. The producer read the same stored
results with its own acceptance program: every target bound to the frozen covers, every outside subbox re-checked in
exact arithmetic against $(3/5+2^{-40})^2$, and $400$ uniformly sampled LG leaves checked at two points each by a
separate exact pointwise computation, without a failure. The lists of shard hashes of the two programs agree in all
nine stages ($15{,}239$ shard files). The two reports and their SHA-256 values are part of the release.

Second-party review (an AI reviewer with separately written programs; this is not peer review): see the table. No
error has been found in the mathematics so far; gaps in the first version of the computation for
Theorem~\ref{thm:L} (certificate data not exported, a floating-point error sum) were found by the review and closed, and a missing argument in the
proof of Theorem~\ref{thm:N} (the winding number of the fan) was found and closed during the review. The reviewer's
program for the box certificates was in turn read line by line by the producer and tested at $2{,}484$ points of
$1{,}222$ of its certificates with a separate exact pointwise computation, without a failure.

\section{Limitations}
Non-convex pentagons are not covered. Numerical experiments suggest that the infimum of $F$ over them is $2\pi^2$,
approached by squares with a vanishing notch; no assertion about general non-convex pentagons is proved here. Both the proof and its independent re-certification
are computer-assisted and were produced by AI systems (Section~\ref{sec:ai}); the mathematical arguments have been
checked by an AI reviewer, not by human experts, and the work has not been peer reviewed.

\section{Roles of AI systems}\label{sec:ai}
Claude (Anthropic) found the methods, wrote the proofs, the producer code and the producer-side checkers, and ran the
computations in a sandbox. GPT (OpenAI) reviewed the mathematics and wrote independent programs for the second-party
checks. The author initiated and coordinated the work, ran the computations on his own computer and decides on the
release. The author did not establish the mathematical correctness by hand; responsibility for the publication lies
with the author.

\appendix
\section{Proofs for Part III}\label{app:far}
Every algebraic identity of this appendix has been verified symbolically, starting from the definitions and
independently of the programs that compute the labels (\texttt{far\_formula\_check.py} in the supplementary material:
$878$ exact checks, the last group comparing the formulas of the programs with the ones derived here).

\begin{lemma}[ears]\label{lem:ear}
Let $P$ be a strictly convex polygon and $C_k>0$. Then the lines of the edges $k-1$ and $k+1$ meet in
$W=V_k+s\,e_{k-1}=V_{k+1}-r\,e_{k+1}$ with $s=B_k/C_k$, $r=A_k/C_k$; the triangle $(V_k,W,V_{k+1})$ lies outside $P$ and has
area $A_kB_k/(2C_k)$; and the polygon $P^+$ obtained by replacing $V_k,V_{k+1}$ by $W$ is strictly convex, with the
edge vectors of $P$ except that $e_k$ is removed and $e_{k\mp1}$ are multiplied by $1+s$, $1+r$.
\end{lemma}
\begin{proof}
$s\,e_{k-1}+r\,e_{k+1}=e_k$ is solved by taking cross products with $e_{k\pm1}$; the area is $\frac12(s\,e_{k-1})\times e_k$,
and $e_k\times(W-V_k)=-sA_k<0$. The consecutive cross products of the new edge vectors are positive, and the edge
directions are those of $P$ without one, so the total turning is $2\pi$.
\end{proof}
A strictly convex pentagon has an edge with $C_k>0$: with the exterior angles $\theta_i\in(0,\pi)$,
$C_k=|e_{k-1}||e_{k+1}|\sin(\theta_k+\theta_{k+1})$, and the five sums $\theta_k+\theta_{k+1}$ add up to $4\pi$.

\paragraph{Normalisation.} $Q=P^+$ for a kept edge $k$ is a strictly convex quadrilateral with the vertex $W$. A
similarity maps a longest edge of $Q$ to $q_0q_1$; the reflection $(x,y)\mapsto(1-x,y)$ exchanges the roles of $q_0,q_1$
and of $q_2,q_3$, so $W\in\{q_3,q_0\}$ after a reflection. The two vertices of $P$ on the edges at $W$ cut away the
fractions $t=s/(1+s)$, $u=r/(1+r)$ (in the appropriate order). If the base is a longest edge then $|q_2-q_1|,|q_3|\le1$,
which gives $q\in\mathfrak B$.

\paragraph{Edge vectors of $P$.} With $e_0,\dots,e_3$ the edge vectors of $Q$:
$j=3$: $(e_0,\ e_1,\ (1-t)e_2,\ te_2+ue_3,\ (1-u)e_3)$; $j=0$: $((1-t)e_0,\ e_1,\ e_2,\ (1-u)e_3,\ te_0+ue_3)$. For the cut
edge ($k=3$ resp.\ $4$) $A=(1-t)u\,\delta$, $B=t(1-u)\delta$, $C=(1-t)(1-u)\delta$ with $\delta=D$ resp.\ $q_{3y}$, so
$\mathrm{ear}_{\rm cut}=tu\,\delta/2$ and $|P|=|Q|(1-tu\tau_j)$.

\paragraph{The comparisons of label X.} For the other four edges one finds
\begin{center}\small\begin{tabular}{@{}cccl@{}}\toprule
$j$ & $k'$ & $C_{k'}$ & $\mathrm{ear}_{k'}<\mathrm{ear}_{\rm cut}$ (given $C_{k'}>0$) is equivalent to\\\midrule
3 & 0 & $(1-u)M$ & $tu>q_{2y}q_{3y}/(DM)$\\
3 & 1 & $(1-t)L$ & $tu>q_{2y}K/(LD)$\\
3 & 2 & $tK-uM$ & $K(1-2t)+tuM<0$\\
3 & 4 & $u\,q_{3y}+t(q_{2y}-q_{3y})$ & $q_{3y}(1-2u)-tu(q_{2y}-q_{3y})<0$\\
0 & 1 & $(1-t)L$ & $tu>q_{2y}K/(q_{3y}L)$\\
0 & 2 & $-(1-u)M$ & $tu>DK/(q_{3y}(-M))$\\
0 & 0 & $uM+t\,q_{2y}$ & $q_{2y}(1-2t)-tuM<0$\\
0 & 3 & $uD+t(q_{2y}-q_{3y})$ & $D(1-2u)-tu(q_{2y}-q_{3y})<0$\\\bottomrule
\end{tabular}\end{center}
For instance $j=3$, $k'=2$: $A=(1-t)K$, $B=(1-t)uD$, $C=tK-uM$, and $(1-t)^2uKD/(2C)<tuD/2$ is $(1-t)^2K<t(tK-uM)$.
On a box, the first kind is tested with $\inf t\cdot\inf u$, the suprema of the factors of the numerator and the
infima of those of the denominator; the second kind at the $64$ corners (the polynomials are multi-affine in each of
the three charts). In the chart E3, with $r=h/q_{2y}$ and $x=q_{2x}$,
$K=q_{2y}(r(x-1)+1-b)$, $L=q_{2y}(r-(1-b))$, $M=q_{2y}(r(x-1)-b)$, $q_{3y}=q_{2y}(b+r)$, $D=q_{2y}\,x\,r$,
so the comparisons of the first kind depend on $(x,b,r,tu)$ only; at the face $q_{2y}=0=h$ they are verified for
$r\in(0,1]$ and, divided by $r^2$, for $1/r\in[0,1]$.

\begin{lemma}[ruled charts]\label{lem:ruled}
In each of the charts std, E2, E3 and for $j=3,0$, every vertex of $P(\zeta)$ is affine in each coordinate, and for
each coordinate $\zeta_m$ the derivatives $\partial_mV_i$, $i=0,\dots,4$, are parallel to one vector (depending on the
other coordinates). Consequently every cross product of two vertex differences, the area, and every linear
functional of the vertices is multi-affine; a multi-affine function $f$ satisfies $f(\zeta)=\sum_c\omega_c(\zeta)f(\zeta_c)$
over the corners $\zeta_c$ of a box, with weights $\omega_c\ge0$, $\sum\omega_c=1$ independent of $f$; and a function that is
concave in each coordinate attains its minimum over a box at a corner.
\end{lemma}
\begin{proof}
The vertices of $P$ are combinations of $q_0,\dots,q_3$ with coefficients $1,t,1-t,u,1-u$. In the chart std a
coordinate of $w$ moves $q_2$ or $q_3$ along a coordinate axis; in E2 the coordinates $q_{3x},q_{3y}$ move $q_3$ along an
axis and $q_2$ by $a$ times the same vector, $a$ moves $q_2$ along $q_3-q_1$ and $h$ moves it along $(0,1)$; in E3
$q_{2x},q_{2y}$ move $q_2$ along an axis and $q_3$ by $b$ times the same, $b$ moves $q_3$ along $q_2$ and $h$ along
$(0,1)$; $t$ and $u$ move one vertex each. The second derivative in $\zeta_m$ of $(V_a-V_b)\times(V_c-V_d)$ is
$2\,\partial_m(V_a-V_b)\times\partial_m(V_c-V_d)=0$. The weights are the products of the one-dimensional barycentric
weights. For the last claim move one coordinate at a time to the end where the function is not larger.
\end{proof}

\paragraph{Hand-over.} Let $q$ be admissible in a std box labelled E2. From $0<K<\frac18(1-q_{3x})$ we get $1-q_{3x}>0$;
with $a=(1-q_{2x})/(1-q_{3x})$ and $h=K/(1-q_{3x})\in(0,\frac18)$ the identity $K=q_{3y}(q_{2x}-1)+q_{2y}(1-q_{3x})$ gives
$q_2=q_1+a(q_3-q_1)+h(0,1)$, and $|a|\,|q_3-q_1|\le|e_1|+h\le\frac98$ with $|q_3-q_1|\ge q_{3y}\ge\frac14$. For E3:
$b=q_{3x}/q_{2x}$, $h=D/q_{2x}\in(0,\frac18)$, $|b|\le|q_{3x}|/q_{2x}\le4$.

\paragraph{S(a).} Choose a non-negative rational $\tau$ such that, at every corner of the box and for every vertex
pair $i,k$, $|P|-\tau|n|^{-2}(p_i-p_k)^2\ge0$, where $p=\langle V,n\rangle$; no division by a corner width is used (a
degenerate corner may have $W_n=0$). For vertices $i,k$ this function is a multi-affine function minus the square of
one, hence concave in each coordinate; by Lemma~\ref{lem:ruled} it is non-negative on the box. So
$|P|\ge\tau|n|^{-2}W_n^2$ on the box. For a strictly convex member $W_n>0$, $P$ lies in a strip of width $W_n/|n|$, and
$F(P)\ge\pi^2|n|^2|P|/W_n^2\ge\pi^2\tau$; the strict label requires $\pi^2\tau>F^*$.

\paragraph{S(b).} For $l>0$ and $\max_kc_k>0$, with $H=\max_kc_k/l$, $P$ lies in a strip of width $H$ and contains a triangle with base $l$ and
height $H$. On a box $\max_kc_k$ is at most the largest corner value, and $l^2$ at least the sum of the squares of the
smallest corner moduli of the two components of $V_b-V_a$ (where these do not change sign).

\paragraph{R.} $P$ is the intersection of its five edge half-planes; those of the edges $m+1,m+3,m+4$ intersect in
$T\supset P$. Lemma~\ref{lem:ear} for the edge $m$ of $P$, then for the edge $m+2$ of $P^+$ (its $A$ and $C$ are those of
$P$ times $1+r$, so the ear is the same) shows that $T$ is a triangle with $|T|=|P|+\mathrm{ear}_m+\mathrm{ear}_{m+2}$.
On a box the ears are bounded by $\sup A\,\sup B/(2\inf C)$ and $|P|$ by its smallest corner value.

\paragraph{O and N.} By Lemma~\ref{lem:ruled} every vertex of $P(\zeta)$ is a convex combination of its corner values,
so $P(\zeta)$ lies in the hull of the $320$ corner vertices, $N_p-r_pD$ has modulus at most its largest corner
modulus, and $|D|\ge\mathrm{Re}(D\bar e)\ge\min_{\rm corners}\mathrm{Re}(D\bar e)$.

\paragraph{Roots.} A cell of $\mathfrak B$ is not a std root only if one of $q_{2y},q_{3y},K,D$ is $\le0$ at its corners
(hence on the cell: no strictly convex $Q$) or one of $|e_1|^2,|e_2|^2,|e_3|^2$ exceeds $1$ on the cell (no admissible
$q$). So every admissible $q$ lies in a root; likewise every strictly convex $Q$ of the edge-chart ranges.

\begin{thebibliography}{99}
\bibitem{ps} G.~P\'olya, G.~Szeg\H{o}, \emph{Isoperimetric Inequalities in Mathematical Physics}, Princeton, 1951.
\bibitem{bb-jep} B.~Bogosel, D.~Bucur, On the polygonal Faber--Krahn inequality, \emph{J. \'Ec. polytech. Math.} 11 (2024) 19--105, doi:10.5802/jep.250.
\bibitem{bb-val} B.~Bogosel, D.~Bucur, Polygonal Faber--Krahn inequality: local minimality via validated computing, arXiv:2406.11575.
\bibitem{wang} H.~Wang, The P\'olya--Szeg\H{o} inequalities for polygons, arXiv:2611.22305 (17 November 2026).
\bibitem{liu} Public repository \texttt{github.com/xfliu/polya-szego-conjecture}, maintained under the GitHub account \texttt{xfliu}; read at commit \texttt{bf7c9c288f56691745f6bcdf31c1128eaa937a93} (September 2026).
\bibitem{cghl} Z.~Cheng, C.~Gui, Y.~Hu, Q.~Li, The P\'olya--Szeg\H{o} conjecture for convex polygons with many sides, arXiv:2609.18500 (2026).
\bibitem{bbf} B.~Bogosel, D.~Bucur, I.~Fragal\`a, Proof of the local version of the P\'olya--Szeg\H{o} conjecture for the torsional rigidity of polygons, arXiv:2609.23051 (2026).
\bibitem{lehmann} N.~J.~Lehmann, Beitr\"age zur numerischen L\"osung linearer Eigenwertprobleme I, II, \emph{Z. Angew. Math. Mech.} 29 (1949) 341--356; 30 (1950) 1--16.
\bibitem{gh} F.~Goerisch, H.~Haunhorst, Eigenwertschranken f\"ur Eigenwertaufgaben mit partiellen Differentialgleichungen, \emph{Z. Angew. Math. Mech.} 65 (1985) 129--135.
\bibitem{bg} H.~Behnke, F.~Goerisch, Inclusions for eigenvalues of selfadjoint problems, in: \emph{Topics in Validated Computations} (Oldenburg, 1993), Stud. Comput. Math. 5, North-Holland, 1994, 277--322.
\bibitem{krahn} E.~Krahn, \"Uber Minimaleigenschaften der Kugel in drei und mehr Dimensionen, \emph{Acta Comm. Univ. Tartu (Dorpat)} A9 (1926) 1--44.
\bibitem{hong} I.~Hong, On an inequality concerning the eigenvalue problem of membrane, \emph{K\={o}dai Math. Sem. Rep.} 6 (1954) 113--114.
\bibitem{flint} The FLINT team, \emph{FLINT: Fast Library for Number Theory}, version 3, \texttt{flintlib.org}; F.~Johansson, Arb: efficient arbitrary-precision midpoint-radius interval arithmetic, \emph{IEEE Trans. Comput.} 66 (2017) 1281--1292.
\bibitem{bbf2} B.~Bogosel, D.~Bucur, I.~Fragal\`a, Proof of the polygonal Saint-Venant inequality conjectured by P\'olya--Szeg\"o, arXiv:2611.24573 (18 November 2026).
\bibitem{v10} P.~Sun, \texttt{pentagon-paqc-proofs}, version~1.0, \texttt{github.com/PYCSON/pentagon-paqc-proofs}, released 1 November 2026.
\bibitem{makai} E.~Makai, On the principal frequency of a membrane and the torsional rigidity of a beam, in: \emph{Studies in Mathematical Analysis and Related Topics}, Stanford Univ. Press, 1962, 227--231.
\end{thebibliography}
\end{document}
